\subsection{From spacing and type counts to distance}

This subsection supplies the analytic bounds used by
Theorem~\ref{thm:baa-fixed}. All limits keep $p$ fixed and take $N=3K$.
For $1\le s\le K$, let $Z_s(d)$ count messages $x\in\F_p^K$ with
$\wt(x)=s$ and $\wt(E_{P_1,P_2}x)\le d$.

\begin{lemma}[Sparse messages]\label{lem:baa-sparse}
Fix $0<\delta<1/2$ and $0<\eta<1$. Put
\[
 c_\delta:=-\tfrac12\log(4\delta(1-\delta)),\quad
 L_\delta:=e^{c_\delta}-1,\quad R_0:=3/\eta,
\]
\[
 g(R):=\frac{(p-1)eR}{3}
             \left(\frac{3}{(R-3)L_\delta}\right)^{3/2}.
\]
If $R_0>6$, $3/((R_0-3)L_\delta)<1$, and $g(R_0)<1$, then
\[
 \sum_{s=1}^{\lfloor\eta K\rfloor}
       \mathbb E Z_s(\lfloor\delta N\rfloor)\longrightarrow0.
\]
\end{lemma}
\begin{proof}
First fix the order of $t$ nonzero increments entering one accumulator.
Their positions under a uniform permutation form a uniform $t$-subset
of $[N]$. Let $a$ be the number of nonzero partial sums after these
increments. The first partial sum is nonzero, and two consecutive partial
sums cannot both vanish. Thus $a\ge\lceil t/2\rceil$.

The gap before the first increment is nonnegative. Each of the $t$ gaps
following an increment is positive, including the final gap through $N$.
These gaps sum to $N$ and are uniform over the $\binom Nt$ possibilities.
If $Y$ is the output weight, counting the gaps at nonzero and zero states gives
\[
 \Pr[Y=w\mid\text{increment order}]
 =\frac{\binom{w-1}{a-1}\binom{N-w}{t-a}}{\binom Nt}.
\]
Summing to $d$ is the probability that a hypergeometric variable with
population $N$, $d$ marked elements, and $t$ draws is at least $a$.
For $\theta>0$, its moment generating function is at most
$(1-d/N+(d/N)e^\theta)^t$. This follows by applying the elementary
symmetric mean inequality to the $N$ weights, each equal to $1$ or
$e^\theta$. The Chernoff bound at threshold $t/2$ therefore gives
\begin{equation}\label{eq:baa-spacing-tail}
 \Pr[Y\le\lfloor\delta N\rfloor]\le e^{-c_\delta t}.
\end{equation}
For any $c>0$, the same gap count and the negative-binomial series give
\begin{equation}\label{eq:baa-spacing-laplace}
 \mathbb E e^{-cY}
 \le \frac{(t)_a}{(N)_a}(e^c-1)^{-a},
 \qquad (v)_a:=v(v-1)\cdots(v-a+1).
\end{equation}
Indeed, replace $\binom{N-w}{t-a}$ by $\binom{N-a}{t-a}$ and use
$\sum_{w\ge a}\binom{w-1}{a-1}e^{-cw}=(e^c-1)^{-a}$.
Both bounds hold uniformly over the nonzero increment values and their order.

A message of support $s$ has $t=3s$ nonzero coordinates after repetition.
Condition on its first accumulator output, of weight $Y$.
Independence of $P_2$ and~\eqref{eq:baa-spacing-tail} bound its final
low-weight probability by $e^{-c_\delta Y}$. For $s\le\eta K$,
applying~\eqref{eq:baa-spacing-laplace} with $c=c_\delta$ then bounds the
probability by
\[
 \left(\frac{t}{(N-t)L_\delta}\right)^{t/2}.
\]
Here $(t)_a/(N)_a\le(t/(N-t))^a$, and the assumed base is less than one.
There are $\binom Ks(p-1)^s\le((p-1)eK/s)^s$ such messages, so
$\mathbb E Z_s(\lfloor\delta N\rfloor)\le g(N/s)^s$.
The function $g(R)$ decreases for $R>3$ and tends to zero as
$R\to\infty$. Each fixed-$s$ term tends to zero, while all terms are
dominated by the summable sequence $g(R_0)^s$. This proves the limit.
\end{proof}

For the remaining messages, let $\mathcal T_p(\eta,\delta)$ be the
stationary types from Appendix~\ref{sec:baa} with
$1-\mu(0)\ge\eta$ and $1-\zeta(0)\le\delta$. Set
\[
 \Lambda_p(\eta,\delta):=
       \max_{\tau\in\mathcal T_p(\eta,\delta)}\Phi_p(\tau).
\]
The maximum of an empty set is $-\infty$.

\begin{lemma}[Bulk messages]\label{lem:baa-bulk}
If $\Lambda_p(\eta,\delta)<0$, then
$\sum_{s=\lceil\eta K\rceil}^{K}
\mathbb E Z_s(\lfloor\delta N\rfloor)\to0$.
\end{lemma}
\begin{proof}
Write $\operatorname{Mult}_\ell(\alpha)$ for the number of length-$\ell$
words of type $\alpha$, and $T_N(\alpha,\beta)$ for the number of
increment words of type $\alpha$ whose accumulator output has type $\beta$.
A permutation is uniform on the words of a fixed type. Thus the expected
number for message, middle, and output types $(\mu,\beta,\zeta)$ is
\[
 \operatorname{Mult}_K(\mu)
 \frac{T_N(\mu,\beta)T_N(\beta,\zeta)}
      {\operatorname{Mult}_N(\mu)\operatorname{Mult}_N(\beta)}.
\]
The first increment type is $\mu$ because repetition preserves its
normalized composition. The zero message is excluded by $\eta>0$.

For each accumulator, record the $p^2$ transition counts $(S,S+D)$.
Given the initial state zero and the ordered successor lists at each
state, the path is determined. The number of paths with a transition
type $Q$ is therefore at most
$\exp(N(H(Q)-H(Q_S)))$, where $Q_S$ is the law of the state $S$
under $Q$.
Use
\[
 (\ell+1)^{-p}e^{\ell H(\alpha)}
 \le\operatorname{Mult}_\ell(\alpha)\le e^{\ell H(\alpha)}.
\]
There are at most $(N+1)^{2p^2+3p}$ choices of the three compositions
and two transition types. Together with the two denominators, the
polynomial prefactor is at most $(N+1)^{2p^2+5p}$.

The source and target state marginals of a path differ only at its
endpoints, by at most $1/N$ in each coordinate. Entropy continuity on
the fixed-dimensional probability simplex and compactness imply
\[
 \limsup_{N\to\infty}\frac1N\log
 \sum_{s=\lceil\eta K\rceil}^{K}
       \mathbb E Z_s(\lfloor\delta N\rfloor)
 \le\Lambda_p(\eta,\delta).
\]
To see uniformity, any sequence of maximizing finite path types has a
convergent subsequence; its limit satisfies the stationary constraints,
and its objective converges to $\Phi_p$ at that limit.
A strictly negative upper bound absorbs the polynomial prefactor.
\end{proof}

\subsection{The certified bulk objectives}

For $p=5,7$, a sign-orbit relaxation bounds the full stationary objective.
Let $[a]:=\{a,-a\}$, let $\mathcal O_p$ be the set of these orbits,
and put
\[
 \mathcal C_p:=\{([d],[s],[t]):d,s,t\in\F_p,\ t=s+d\}.
\]
Optimize over laws $\gamma_1,\gamma_2$ supported on $\mathcal C_p$.
Their increment, source, and target marginals must respectively be
$(\bar\mu,\bar\beta,\bar\beta)$ and
$(\bar\beta,\bar\zeta,\bar\zeta)$.
Let $\pi_i$ be the increment--source marginal of $\gamma_i$. The objective is
\begin{equation}\label{eq:baa-orbit-objective}
 \widehat\Phi_p:=H(\pi_1)+H(\pi_2)-\tfrac23H(\bar\mu)
              -2H(\bar\beta)-H(\bar\zeta)
              +\tfrac13(1-\bar\mu([0]))\log2.
\end{equation}
Every full type projects to a feasible orbit type with the same support
weights. Data processing gives $I(D;S)\ge I([D];[S])$, and
$H(\mu)\le H(\bar\mu)+(1-\bar\mu([0]))\log2$.
Since $\Phi_p=H(\mu)/3-I_{Q_1}(D;S)-I_{Q_2}(D;S)$, these inequalities
prove $\Phi_p\le\widehat\Phi_p$.

Over $\F_3$, a sharper symmetry reduction is exact. For a distribution
$\rho$, let $\sigma$ negate every symbol and define
$\operatorname{Asym}(\rho):=H((\rho+\sigma\rho)/2)-H(\rho)$.

\begin{lemma}[Ternary sign reduction]\label{lem:baa-ternary-sign}
A stationary transition law $T$ on $\F_3^2$, with state marginal
$\beta$ and increment marginal $\alpha$, satisfies
\[
 \operatorname{Asym}(T)\ge
       \operatorname{Asym}(\alpha)+\operatorname{Asym}(\beta).
\]
Consequently, the maximum defining $\Lambda_3(\eta,\delta)$ is attained
by sign-symmetric compositions and transition laws whenever its domain is nonempty.
\end{lemma}
\begin{proof}
For nonnegative $u,v$, put
$j(u,v):=u\log(2u/(u+v))+v\log(2v/(u+v))$, with continuous boundary
values. This function is convex and positively homogeneous.
The off-diagonal entries of a stationary three-state law form a circulation.
After exchanging symbols 1 and 2 if necessary, write
\[
 T_{01}=c+q,\ T_{10}=c,\ T_{12}=h+q,\ T_{21}=h,
 \ T_{20}=b+q,\ T_{02}=b,
\]
with $b,c,h,q\ge0$, and put $T_{11}=y$, $T_{22}=z$.
The three asymmetries are
\begin{align*}
 \operatorname{Asym}(T)&=j(c+q,b)+j(c,b+q)+j(y,z)+j(h+q,h),\\
 \operatorname{Asym}(\alpha)&=j(b+c+h+3q,b+c+h),\\
 \operatorname{Asym}(\beta)&=j(c+y+h+q,b+z+h+q).
\end{align*}
Subadditivity of $j$ bounds the last expression by
$j(y,z)+j(c+h+q,b+h+q)$.
Set $m=b+c$, $x=c-b$, and $J(s,d):=j((s+d)/2,(s-d)/2)$.
It suffices to prove that
\begin{align*}
 G(x):={}&J(m+q,x+q)+J(m+q,x-q)+J(2h+q,q)\\
        &-J(2m+2h+3q,3q)-J(m+2h+2q,x)
 \ \ge0.
\end{align*}
This is an even convex function of $x$. Indeed,
$\partial_d^2J(s,d)=s/(s^2-d^2)$ is convex in $d$ and decreases in $s$;
hence the first two terms' second derivatives dominate the last one's.
The feasible interval $[-m,m]$ is symmetric, so the minimum is at zero.
There, subadditivity of $J$ gives
\[
 G(0)=2J(m+q,q)+J(2h+q,q)-J(2m+2h+3q,3q)\ge0.
\]
Continuity covers boundary points. Finally average both stationary laws
and their compositions with their sign images. The change in $\Phi_3$ is
\[
 \operatorname{Asym}(T_1)+\operatorname{Asym}(T_2)
 -\tfrac23\operatorname{Asym}(\mu)-2\operatorname{Asym}(\beta)
 -\operatorname{Asym}(\zeta)\ge\tfrac13\operatorname{Asym}(\mu)\ge0.
\]
The support constraints are preserved.
\end{proof}

A symmetric transition with nonzero increment fraction $\nu$ and nonzero
state fraction $\nu'$ has matrix
\[
 T(\nu,\nu',b)=
 \begin{pmatrix}
 1-\nu'-2b&b&b\\
 b&(\nu'-\nu+2b)/2&\nu/2-2b\\
 b&\nu/2-2b&(\nu'-\nu+2b)/2
 \end{pmatrix},
\]
where
\[
 \max(0,(\nu-\nu')/2)\le b\le\min((1-\nu')/2,\nu/4).
\]
Let $\nu_1,\nu_2,\nu_3$ be the nonzero fractions of $\mu,\beta,\zeta$.
Thus the exact ternary bulk objective maximizes
$\Phi_3$ using $T(\nu_1,\nu_2,b_1)$ and $T(\nu_2,\nu_3,b_2)$ over
$\eta\le\nu_1\le1$, $\nu_1\le2\nu_2$, $\nu_2\le2\nu_3$, and
$0\le\nu_3\le\delta$.
For each $(\nu,\nu')$ the entropy is concave in $b$; an interior maximizer
solves
\[
 (1-\nu'-2b)(\nu/2-2b)^2=b^2(\nu'-\nu+2b)/2.
\]
This specifies the objective checked by the ternary verifier, including
its feasible boundary.

The directed-arithmetic verifiers give the following global upper bounds.
The displayed values are rounded upward.
\begin{center}
\begin{tabular}{@{}cccl@{}}
\toprule
$p$ & $\eta$ & $\delta$ & upper bound and objective \\
\midrule
3 & $1/10000$ & $49/250$ & $-3.6177\cdot10^{-7}$, exact symmetric \\
5 & $1/60$ & $1/20$ & $-1.1386\cdot10^{-6}$, sign orbit \\
7 & $1/16$ & $9/500$ & $-1.1596\cdot10^{-4}$, sign orbit \\
\bottomrule
\end{tabular}
\end{center}
The orbit verifier covers the feasible composition domain by rational boxes.
It upper-bounds positive joint entropies by cross-entropy at positive trial
laws, and negative marginal entropies by secants of $x\log x$.
The resulting linear program includes the orbit incidence and marginal
constraints. A rational dual solution bounds its objective on the whole box;
directed logarithm intervals make each bound outward. Boxes discarded as
infeasible have exact dual witnesses. Multiplicative automorphisms supply
additional valid cuts for $p=7$. The ternary verifier instead bounds the
concave transition entropies by supporting planes and the negative marginal
entropies by their endpoint secants. Rational boxes and interval arithmetic
enclose the maximizing values of $b_1,b_2$ and the resulting objective.

The certificates and their verifiers are listed in
Appendix~\ref{app:reproduction}.

\begin{proof}[Proof of Theorem~\ref{thm:baa-fixed}]
Apply Lemma~\ref{lem:baa-sparse} with the three $(p,\eta,\delta)$ rows
above. For $p=3$, $L_\delta\ge1/4$ and $R_0=30000$, giving
$g(R_0)<1/2$ from
$16\cdot30000^2\cdot12^3<29997^3$ and $e<3$.
For $p=5$, $L_\delta>9/7$ and $R_0=180$ give
$g(R_0)<660(7/531)^{3/2}<1$, using $e<11/4$ and
$660^2\cdot7^3<531^3$.
For $p=7$, $L_\delta>367/133$ and $R_0=48$ give
$g(R_0)<264(133/5505)^{3/2}<1$, using $e<11/4$ and
$264^2\cdot133^3<5505^3$. The other sparse hypotheses hold in all three rows.
The certified bulk bounds and Lemma~\ref{lem:baa-bulk} handle every larger
support. Markov's inequality on the total nonzero-message count proves
the three distance claims.
\end{proof}

\begin{proof}[Proof of Proposition~\ref{prop:sparse-quinary}]
Take $p=5$, $\delta=131/500=0.262$, $\eta=10^{-5}$, and $R_0=300000$ in
Lemma~\ref{lem:baa-sparse}.
The comparison displayed after Proposition~\ref{prop:sparse-quinary} gives
$L_\delta>3/22$.
Then $g(R_0)<4R_0(22/(R_0-3))^{3/2}<4/5$, using $e<3$ and
$25\cdot300000^2\cdot22^3<299997^3$.  The other hypotheses of the lemma
hold, and Lemma~\ref{lem:baa-sparse} proves the stated vanishing expectation.
\end{proof}
